\section{Background}

This section considers three uses of time in puzzle games: changing temporal rules, learning a recurring schedule, and carrying knowledge from one attempt to the next. It then reviews action replay, state management, and level authoring as design and implementation concerns related to this project.

\subsection{Time-Loop and Time-Manipulation Games}

\textit{Braid} makes rewind part of its puzzle rules. Players use it to alter the relationship between movement and events in the level; some worlds also introduce objects or situations that do not follow the ordinary rewind rule. Solving a puzzle therefore requires reasoning about both position and which events can be reversed. Stamenkovi\'c and Ja\'cevi\'c analyze the game through the relations among time, space, motion, and causality.\footnote{Du\v{s}an Stamenkovi\'c and Milan Ja\'cevi\'c, \textit{Time, Space, and Motion in Braid: A Cognitive-Semantic Approach to a Video Game}, \textit{Games and Culture} 10, no. 2 (2015): 178--203, \url{https://doi.org/10.1177/1555412014557640}.} Jonathan Blow's GDC presentation on \textit{Braid}'s rewind system also describes the engineering requirement to reproduce past game states accurately while keeping a long rewind history within the memory limits of consoles.\footnote{Jonathan Blow, \textit{The Implementation of Rewind in Braid}, Game Developers Conference, 2010, \url{https://www.gdcvault.com/play/1012850/The-Implementation-of-Rewind-in}.} These works make \textit{Braid} relevant both as a puzzle-design example and as a technical case in temporal state restoration.

\textit{The Sexy Brutale} uses a different structure. The events in its mansion recur on a shared daily schedule, and the player observes characters' routines before intervening to prevent their deaths. The developers describe the characters as following synchronized schedules across the mansion, with events in one location affecting what others hear or do elsewhere.\footnote{Cavalier Game Studios, \textit{Murder Mystery Game The Sexy Brutale Available Now on Xbox One}, Xbox Wire, April 12, 2017, \url{https://news.xbox.com/en-us/2017/04/12/the-sexy-brutale-available-now-xbox-one/}.} Its central puzzle activity is to learn the timeline and identify where an intervention can change its outcome. The repeated setting makes observations useful: the player can compare what happens at the same point in successive cycles.

\textit{Twelve Minutes} places the emphasis on knowledge retained by the player. Director Luis Antonio described accumulated knowledge as the concept behind the game and explained that short loops and a small environment help players understand the consequences of their actions and plan another attempt.\footnote{Bryant Francis, \textit{Using Accumulated Player Knowledge to Drive Design Decisions in Twelve Minutes}, interview with Luis Antonio, \textit{Game Developer}, September 9, 2021, \url{https://www.gamedeveloper.com/design/q-a-i-twelve-minutes-i-director-luis-antonio-on-making-good-time-loop-games}.} The character and situation repeat, but the player uses what they have learned to choose different actions. This is relevant to the project's compact-room format and its aim of making attempts quick to inspect and revise.

\subsection{Puzzle Game Design}

Puzzle design depends on players being able to understand the goal, the available actions, and the consequences of those actions. Schell and Adams discuss goals, rules, challenge, and feedback as core parts of game design.\footnote{Jesse Schell, \textit{The Art of Game Design: A Book of Lenses}, 3rd ed. (CRC Press, 2019); Ernest Adams, \textit{Fundamentals of Game Design} (New Riders, 2014).} Applied to this project, the relevant principle is to introduce mechanics in a form players can observe, then combine familiar rules to create more demanding puzzles. A room should provide enough feedback for the player to understand what a ghost did and why an interaction succeeded or failed. Without that feedback, replay is difficult to reason about.

\subsection{Action Recording and Replay}

\textit{The Misadventures of P.B. Winterbottom} provides a close example of recorded behaviour becoming part of a puzzle solution: the player records actions and cooperates with earlier versions of the character.\footnote{\textit{The Misadventures of P.B. Winterbottom}, Steam store description, \url{https://store.steampowered.com/app/40930/The_Misadventures_of_PB_Winterbottom/}.} This differs from a passive replay or visual trail because the recorded character can interact with the level alongside the current one.

The proposed game uses a related mechanic in turn-based rooms. Its design records player commands so a ghost can repeat the decisions that produced a useful position or interaction. The room rules then determine the outcome of those commands during replay. This means that replay must specify what happens when the room has changed since the original attempt---for example, when a move is blocked or an interaction target is no longer available. Defining those cases is part of making the ghost's behaviour predictable.

\subsection{State Management}

A loop reset has to restore the room without erasing the history needed by its ghosts. The level's initial object positions and mechanism states are resettable; completed action sequences persist across attempts. A full level restart has a wider scope and clears those histories as well. Keeping these categories distinct gives the game a clear basis for resetting the room and replaying earlier attempts.

Statecharts provide a formal way to describe system states and transitions, while Nystrom's programming-pattern examples discuss ways of organizing object behaviour and state in software.\footnote{David Harel, \textit{Statecharts: A Visual Formalism for Complex Systems}, \textit{Science} 231, no. 4737 (1987): 671--676; Robert Nystrom, \textit{Game Programming Patterns} (Genever Benning, 2014).} For this project, the practical implication is to define the order of a turn explicitly: collect an action, resolve its effects and environmental responses, then advance the turn. This helps prevent the result from depending on incidental update order.

\subsection{Level Representation and Creation}

A data-driven level separates puzzle content from the code that runs it. In this project, level data is intended to describe the room layout, starting conditions, objectives, turn allowance, and interactive objects; JSON provides the save and load format for the level workflow. Although procedural-content research also studies representations for game content, the project uses designed rooms rather than automatic level generation. The related concern is the need for a representation that the authoring workflow and game runtime can both interpret consistently.\footnote{Noor Shaker, Julian Togelius, and Mark J. Nelson, \textit{Procedural Content Generation in Games} (Springer, 2016).}

\subsection{Comparison with the Proposed Game}

These games provide distinct precedents rather than a single equivalent system. \textit{Braid} makes temporal manipulation part of the puzzle rules; \textit{The Sexy Brutale} asks players to learn and alter a recurring event schedule; and \textit{Twelve Minutes} foregrounds knowledge accumulated across repeated attempts. \textit{The Misadventures of P.B. Winterbottom} is a closer precedent for replayed actions cooperating with the current character.

The proposed game brings these ideas into a compact, turn-based room: each attempt has a limited action allowance, and past commands are replayed by ghosts that share the room with the player. The key distinction is that earlier attempts are not only remembered or used to alter time; they remain as actors whose actions can affect the current puzzle. This combination motivates the project's focus on consistent replay, clearly visible interactions, and a reliable separation between resettable room state and persistent action histories.
